
\Problem{Encuestas UPC}

\Statement

La UPC necesita un programa para procesar las encuestas de satisfacción que responden los estudiantes. El programa debe leer a la entrada todas las encuestas de todas las asignaturas y producir un listado ordenado de los resultados para cada asignatura.

\medskip

El número $n$ de preguntas para cada encuesta se indica al principio, con $3 \leq n \leq 20$. Cada encuesta $e$ tiene el código de la asignatura y las respuestas a las $n$ preguntas de la encuesta, que son valores numéricos $e_i$, $1 \leq i \leq n$ con $0 \leq e_i \leq 5$. El resultado para cada asignatura es la lista de las medias de los valores de cada pregunta de la encuesta.

\medskip

Para calcular la media de forma más fiable estadísticamente, no se tienen en cuenta para el cálculo la encuesta con mejor resultado y la encuesta con peor resultado, porque se asume que ambos extremos son encuestas con sesgo. Solo se quita una sola encuesta de cada extremo, aunque haya más de una encuesta con peor o mejor resultado. No obstante, esta corrección solo se aplica cuando la cantidad de encuestas disponible para una asignatura es mayor o igual a 5, para remediar la escasez de muestras.

\medskip

\textbf{Comparación de dos encuestas}: una encuesta $a$ con valores $a_i$ es mayor que otra encuesta $b$ con valores $b_i$ si la suma $s_a = \sum_{i} a_i$ es mayor que la suma $s_b = \sum_{i} b_i$.

\Input

La entrada empieza con un natural $n$, el número de preguntas de cada encuesta. Luego viene un natural positivo $m$, que es el número total de encuestas. A continuación vienen las encuestas, cada una con un código (un natural) y los valores $e_i$, $1 \leq i \leq n$, con $1 \leq e_i \leq 5$. Las encuestas tienen un orden aleatorio.

\Output

La salida es una lista de los resultados, ordenados por código de asignatura, donde cada resultado contiene el código de la asignatura, seguido por las $n$ medias de cada pregunta de la encuesta. Los resultados deben tener exactamente dos decimales.

\Observation

Recordada que podeis escribir un real \texttt{x} con dos decimales con el comando \texttt{print(f"\{x:.2f\}")}.

\Sample
